\documentclass[10pt,a4paper]{amsart}
\usepackage[margin=19mm]{geometry}
\usepackage{amsmath,amssymb,mathtools}
\usepackage[scr=rsfs,cal=euler]{mathalfa}
\usepackage{xfrac}
\usepackage[unicode,hidelinks,bookmarks=false]{hyperref}
\newcommand{\ie}{i.e.\ }
\newcommand{\cf}{cf.\ }
\newcommand{\resp}{respectively\ }
\newcommand{\Subsection}{\subsection}
\providecommand{\nobreakdash}{\nobreak\hbox{-}}
\setlength{\emergencystretch}{2em}
\multlinegap0pt
\usepackage{bm}
\usepackage{amssymb}
\usepackage{mathtools}
\usepackage[normalem]{ulem}
\usepackage{tikz}
\usepackage{algorithm}\usepackage{algpseudocode}
\usepackage{xcolor}
\newtheorem{assumption}{Assumption}
\newtheorem{lemma}{Lemma}
\newcommand{\Y}[1][]{%
  \ifx\relax#1\relax
    \boldsymbol{Y}_{t-1, 1, s}^{i}%
  \else
    \boldsymbol{Y}_{t-1, 1, s_{#1}}^{i_{#1}}%
  \fi
}
\title{On the Coalescence Time Distribution in Multi-type Supercritical Branching Processes}
\author{Janique Krasnowska, Paul Jenkins, Adam Johansen}
\date{Source: arXiv:2603.11990v1 (CC BY 4.0)}
\begin{document}
\maketitle

\noindent\emph{Source and licence.} On the Coalescence Time Distribution in Multi-type Supercritical Branching Processes. Version arXiv:2603.11990v1 (CC BY 4.0), distributed by arXiv under the Creative Commons Attribution 4.0 International licence, http://creativecommons.org/licenses/by/4.0/, as recorded in the arXiv licence metadata for this exact version and shown on the abstract page https://arxiv.org/abs/2603.11990v1. Full licence notice: \texttt{licenses/arxiv-2603.11990-CC-BY-4.0.txt}.\par\smallskip

\noindent\textit{Editorial scope.} Counted unit: the article's lemma y\_upper (source0.tex lines 1121--1126). The article's complete original proof of it is delivered with it (source0.tex lines 1127--1183, one \texttt{proof} environment of 57 lines). No mathematics is written, repaired or re-derived: the statement and the proof are the article's own lines, in the article's own notation, and no retained line is substituted or re-flowed.  Retained uncounted as the source-local context the proof needs: lines 189--191; lines 198--200; lines 222--224.  Nothing was omitted from any retained span, so the package carries no omission record.  No retained line contains a citation, so the wrapper carries no bibliography and no bibliography entry.  No retained line contains a figure environment, an image-inclusion command or drawing code, so no image is delivered and the package declares no asset dependency.  Editorial formatting only, in addition: an AMS \texttt{amsart} preamble that restates the article's own declarations and macros used by the retained lines, namely the environments \texttt{assumption}, \texttt{lemma} on separate counters, together with the standard packages those lines need.  The retained environments print in this document as Assumption 1, Lemma 1 in order of appearance, whereas the article prints the same items with the numbering of its own full document.  The complete native source of the article as distributed in the arXiv e-print for this exact version is bundled unchanged as \texttt{sources/196249/source0.tex}.
\begin{equation}\label{eq:Zdecomposition}
\boldsymbol{Z}^{i_0}_T = \sum_{i\in S}\sum_{s=1}^{Z_t^{i_0, i}} \boldsymbol{Z}^i_{t,T-t,s},
\end{equation}

\begin{assumption}\label{infinite_matrix}
    $\boldsymbol{M}$ is irreducible and all $m_{ij}^{(r)}$, $r \in \mathbb{N}_+$, $i,j \in S$, are finite.
\end{assumption}

\begin{assumption}\label{R-positive}
  $M_{ij}(R) = \infty$ for any pair $i,j \in S$, $\boldsymbol{u} \cdot\boldsymbol{\nu} < \infty$, and $R \in (0, 1)$ (the process is supercritical).
\end{assumption}

\begin{lemma}\label{y_upper}
Under Assumptions \ref{infinite_matrix} and \ref{R-positive}, for $r\in\mathbb{N}_+$, %\paul{Really not allowing $r=1$ here?}
    \begin{equation*}
       \mathbb{E}\left(\frac{1}{|\boldsymbol{Y}_t|^r}\right) \le   \left(\mathbb{E}\left(\sup_{i \in S}\frac{1}{|\boldsymbol{Y}_{1}^{i}|}\right)\right)^{tr} .
    \end{equation*}
\end{lemma}

\begin{proof}
Denote by $\boldsymbol{Y}_{t-1, 1, s}^i$ the transformed branching process started from an ancestor of type $i$ with index $s$ in generation $t-1$, and allowed to evolve for one generation. Using the tower rule, let us condition on the population size at generation $t-1$. For $r \in \mathbb{N}_+$,
\begin{align}
    \mathbb{E}\left(\frac{1}{|\boldsymbol{Y}_t|^r}\right) %\nonumber \\
    %&= \mathbb{E}\left(\mathbb{E}\left(\frac{|\boldsymbol{Y}_{t-1}|^r}{|\boldsymbol{Y}_{t-1}|^r}\frac{1}{|\boldsymbol{Y}_t|^r}\bigm| \boldsymbol{Y}_{t-1}\right)\right) \nonumber \\
    &= \mathbb{E}\left(\frac{1}{|\boldsymbol{Y}_{t-1}|^r}\mathbb{E}\left(\frac{|\boldsymbol{Y}_{t-1}|^r}{|\boldsymbol{Y}_t|^r}\bigm| \boldsymbol{Y}_{t-1}\right)\right) \nonumber \\
    &= \mathbb{E}\left(\frac{1}{|\boldsymbol{Y}_{t-1}|^r}\mathbb{E}\left(\frac{|\boldsymbol{Y}_{t-1}|^r}{(\sum_{i \in S}\sum_{s=1}^{Y_{t-1}^i}|\Y|)^r}\bigm| \boldsymbol{Y}_{t-1}\right)\right) \label{eq:arithmetic_mean} \\
    &= \mathbb{E}\left(\frac{1}{|\boldsymbol{Y}_{t-1}|^r}\mathbb{E}\left(\frac{|\boldsymbol{Y}_{t-1}|^r}{\sum_{i_1 \in S}\sum_{s_1=1}^{Y_{t-1}^{i_1}} \dots \sum_{i_r \in S}\sum_{s_r=1}^{Y_{t-1}^{i_r}}\prod_{j=1}^r|\Y[j]|}\bigm| \boldsymbol{Y}_{t-1}\right)\right), \label{eq:expanded} 
\end{align}
where \eqref{eq:arithmetic_mean} uses a decomposition analogous to \eqref{eq:Zdecomposition}, i.e.\ 
%comes from expressing the full population size in terms of families started at generation $t-1$ and ran for one generation, i.e. 
$
\boldsymbol{Y}_t = \sum_{i \in S}\sum_{s=1}^{Y_{t-1}^i}\Y
$.
%and \eqref{eq:expanded} expands the $r$th power of $\sum_{i \in S}\sum_{s=1}^{Y_{t-1}^i}|\Y|$.
The expression in the inner expectation in \eqref{eq:expanded} is the reciprocal of the arithmetic mean of the collection of $\prod_{j=1}^r|\Y[j]\vert$.
%\begin{equation*}
%    \frac{\sum_{i_1 \in S}\sum_{s_1=1}^{Y_{t-1}^{i_1}} \dots \sum_{i_r \in S}\sum_{s_r=1}^{Y_{t-1}^{i_r}}\prod_{j=1}^r|\Y[j]|}{|\boldsymbol{Y}_{t-1}|^r}
%\end{equation*}
We will use the fact that the arithmetic mean is always greater than or equal to the harmonic mean (and hence the inverse of harmonic mean is greater than or equal to the inverse of the arithmetic mean);
here
\begin{equation}
   \frac{|\boldsymbol{Y}_{t-1}|^r}{\sum_{i_1 \in S}\sum_{s_1=1}^{Y_{t-1}^{i_1}} \dots \sum_{i_r \in S}\sum_{s_r=1}^{Y_{t-1}^{i_r}}\prod_{j=1}^r|\Y[j]|} 
   \le \frac{\sum_{i_1 \in S}\sum_{s_1=1}^{Y_{t-1}^{i_1}} \dots \sum_{i_r \in S}\sum_{s_r=1}^{Y_{t-1}^{i_r}}\frac{1}{\prod_{j=1}^r|\Y[j]|}}{|\boldsymbol{Y}_{t-1}|^r}. \label{eq:harmonic_arithmetic}
\end{equation}
The harmonic mean is defined since we know $|\Y[j]| > 0 $ for any $s_j \in \mathbb{N}_+$ and $i_j \in S$ as by construction of the Harris--Savastyanov transformation the extinction probability is 0.

Let us use inequality \eqref{eq:harmonic_arithmetic} to bound \eqref{eq:expanded}:
\begin{align}
\mathbb{E}\left(\frac{1}{|\boldsymbol{Y}_t|^r}\right)
 %&=\mathbb{E}\left(\frac{1}{|\boldsymbol{Y}_{t-1}|^r}\mathbb{E}\left(\frac{|\boldsymbol{Y}_{t-1}|^r}{\sum_{i_1 \in S}\sum_{s_1=1}^{Y_{t-1}^{i_1}} \dots \sum_{i_r \in S}\sum_{s_r=1}^{Y_{t-1}^{i_r}}\prod_{j=1}^r|\Y[j]|}| \boldsymbol{Y}_{t-1}\right)\right) \nonumber \\
&\le \mathbb{E}\left(\frac{1}{|\boldsymbol{Y}_{t-1}|^r}\mathbb{E}\left(\frac{\sum_{i_1 \in S}\sum_{s_1=1}^{Y_{t-1}^{i_1}} \dots \sum_{i_r \in S}\sum_{s_r=1}^{Y_{t-1}^{i_r}}\frac{1}{\prod_{j=1}^r|\Y[j]|}}{|\boldsymbol{Y}_{t-1}|^r}| \boldsymbol{Y}_{t-1}\right)\right) \nonumber \\
& =\mathbb{E}\left(|\boldsymbol{Y}_{t-1}|^{-2r}\sum_{i_1 \in S}\sum_{s_1=1}^{Y_{t-1}^{i_1}} \dots \sum_{i_r \in S}\sum_{s_r=1}^{Y_{t-1}^{i_r}}\mathbb{E}\left(\frac{1}{\prod_{j=1}^r|\Y[j]|} \mid \boldsymbol{Y}_{t-1}\right)\right) \label{eq:tonelli}\\ 
& =\mathbb{E}\left(|\boldsymbol{Y}_{t-1}|^{-2r}\sum_{i_1 \in S}\sum_{s_1=1}^{Y_{t-1}^{i_1}} \dots \sum_{i_r \in S}\sum_{s_r=1}^{Y_{t-1}^{i_r}}\prod_{j=1}^r\mathbb{E}\left(|\Y[j]|^{-1} \mid \boldsymbol{Y}_{t-1}\right)\right)\label{eq:independenceY}
\end{align}%
where \eqref{eq:tonelli} follows by Tonelli's Theorem and \eqref{eq:independenceY} is due to the independence of the branching processes started from each individual at generation $t-1$. Observe that the branching processes started at generation $t-1$, $\{\Y[j]\}_{j \in [r]}$, are independent of the population size at time $t-1$ and thus the conditional expectation in \eqref{eq:independenceY} is equal to the unconditional expectation:
\begin{align}
&\mathbb{E}\left(|\boldsymbol{Y}_{t-1}|^{-2r}\sum_{i_1 \in S}\sum_{s_1=1}^{Y_{t-1}^{i_1}} \dots \sum_{i_r \in S}\sum_{s_r=1}^{Y_{t-1}^{i_r}}\prod_{j=1}^r\mathbb{E}\left(|\Y[j]|^{-1} \mid \boldsymbol{Y}_{t-1}\right)\right) \nonumber\\ 
& =\mathbb{E}\left(|\boldsymbol{Y}_{t-1}|^{-2r}\sum_{i_1 \in S}\sum_{s_1=1}^{Y_{t-1}^{i_1}} \dots \sum_{i_r \in S}\sum_{s_r=1}^{Y_{t-1}^{i_r}}\prod_{j=1}^r\mathbb{E}\left(|\boldsymbol{Y}_{t-1,1,1}^{i_j}|^{-1}\right)\right) \label{eq:drop_index}\\
& \le\mathbb{E}\left(|\boldsymbol{Y}_{t-1}|^{-2r}\sum_{i_1 \in S}\sum_{s_1=1}^{Y_{t-1}^{i_1}} \dots \sum_{i_r \in S}\sum_{s_r=1}^{Y_{t-1}^{i_r}}\left(\mathbb{E}\left(\sup_{j \in [r]}|\boldsymbol{Y}_{t-1,1,1}^{i_j}|^{-1}\right)\right)^r\right), \label{eq:sup} 
\end{align}
where \eqref{eq:drop_index} follows from the fact that for a given type $i_j$, $\mathbb{E}(|\Y[j]|^{-1}) = \mathbb{E}(|\boldsymbol{Y}_{t-1,1,1}^{i_j}|^{-1})$ and \eqref{eq:sup} upper bounds each element of the product by its supremum over all types.

Note that the supremum over types $\{i_j\}_{j=1}^r$ in \eqref{eq:sup} can be upper bounded by the supremum over \emph{all} possible types:
\begin{align}
&\mathbb{E}\left(|\boldsymbol{Y}_{t-1}|^{-2r}\sum_{i_1 \in S}\sum_{s_1=1}^{Y_{t-1}^{i_1}} \dots \sum_{i_r \in S}\sum_{s_r=1}^{Y_{t-1}^{i_r}}\left(\mathbb{E}\left(\sup_{j \in [r]}|\boldsymbol{Y}_{t-1,1,1}^{i_j}|^{-1}\right)\right)^r\right) \nonumber\\
    & \le  \mathbb{E}\left(|\boldsymbol{Y}_{t-1}|^{-2r}\sum_{i_1 \in S}\sum_{s_1=1}^{Y_{t-1}^{i_1}} \dots \sum_{i_r \in S}\sum_{s_r=1}^{Y_{t-1}^{i_r}}\left(\mathbb{E}\left(\sup_{j \in S}|\boldsymbol{Y}^j_{t-1, 1, 1}|^{-1}\right)\right)^r\right)  \nonumber \\
    & = \mathbb{E}\left(|\boldsymbol{Y}_{t-1}|^{-r}\left(\mathbb{E}\left(\sup_{j \in S}|\boldsymbol{Y}^j_{t-1, 1, 1}|^{-1}\right)\right)^r\right) \nonumber\\
    & = \mathbb{E}\left(|\boldsymbol{Y}_{t-1}|^{-r}\right)\left(\mathbb{E}\left(\sup_{j \in S}|\boldsymbol{Y}^j_{t-1, 1, 1}|^{-1}\right)\right)^r = \mathbb{E}\left(|\boldsymbol{Y}_{t-1}|^{-r}\right)\left(\mathbb{E}\left(\sup_{j \in S}|\boldsymbol{Y}_1^j|^{-1}\right)\right)^r  \label{eq:split2}%\label{eq:gent1_vs1}
    \end{align}
where the first equality in \eqref{eq:split2} follows by the independence of the two quantities and the second is due to the equality in law of $|\boldsymbol{Y}^j_{t-1, 1, 1}|^{-1}$ and $|\boldsymbol{Y}_1^j|^{-1}$.

Iterating this procedure over $t$, we obtain the desired inequality. %get an upper bound for the $r$th harmonic moment of $|\boldsymbol{Y}_t|$ that depends only upon the properties of the process at generation $1$ 
%\begin{equation*}
%    \mathbb{E}\left(\frac{1}{|\boldsymbol{Y}_t|^r}\right) \le   \left(\mathbb{E}\left(\sup_{i \in S}\frac{1}{|\boldsymbol{Y}_{1}^{i}|^r}\right)\right)^t.
%\end{equation*}
\end{proof}

\end{document}
